\chapter{Mélangeur de couleurs}
\section{Sélectionner une famille chromatique}
Les huit cercles sélectionnent les rouges, oranges, jaunes, verts, cyans,
bleus, violets ou magentas. Violet et magenta sont deux bandes distinctes.
Le nom coloré confirme la bande active. Changer de cercle change les
contrôles affichés, sans effacer les réglages des autres bandes.

\section{Teinte, saturation et luminance}
Chaque bande dispose de trois contrôles de $-100$ à $+100$, référence zéro.
La teinte déplace la famille vers ses voisines, la saturation modifie sa
richesse colorée et la luminance son niveau lumineux. Les familles se
recouvrent progressivement : une couleur située à une frontière peut
réagir à plusieurs bandes. Il ne s’agit pas d’une sélection d’objet.

\paragraph{Exemple.}
Pour travailler une fleur orange, sélectionner d’abord les oranges,
essayer une faible modification de luminance puis de saturation et examiner
les rouges et jaunes voisins. Si le rendu de départ est déjà trop rouge,
vérifier d’abord profil et balance des blancs plutôt que compenser systématiquement
ce défaut avec le mélangeur.

Comparer avant/après avec \touche{Ctrl + Maj + Y}. Aucun raccourci de sélection
des huit cercles n’est déclaré ; utiliser les cercles ou la navigation de focus.
\gestesreglages
